\section{Conclusions}

We have examined the metabolic cost of maintaining body shape against gravity across a 12-row vertebrate panel and report, as a negative result, that the Bio-Gravitational Number $\mathcal{B}_g = EI/MgL^2$ does not isolate the human condition: no threshold exists whose crossing predicts scoliosis. Its critical value depends on boundary conditions, load distribution and support, so it is a descriptor of passive gravitational resistance rather than a universal criterion. In this model, adolescent idiopathic scoliosis is instead a possible consequence of a rate competition during growth --- growth-driven remodeling outpacing geometric restoration --- offered as a falsifiable hypothesis, not a demonstrated cause (extension to bipedal and facultatively bipedal species beyond hominins remains an explicit next step). Exploratory AlphaFold structural analysis generates hypotheses regarding the molecular architecture of this vulnerability, suggesting a potential demand-supply structural gap between mechanosensory and metabolic proteins. The recovery-ratchet framework unifies growth biomechanics, active geometric restoration and developmental metabolism under a single rate competition, generating three falsifiable predictions testable with existing technology; the delay-Hopf alternative was tested and rejected (Supplementary Section~\ref{sec:supp_hopf_rejected}). These results motivate a research agenda in which the metabolic supply available during the adolescent growth spurt is examined as a candidate modifiable risk factor for AIS progression. We do not propose, on the basis of in-silico modeling alone, any specific clinical intervention; the framework's value at this stage is in identifying candidate mechanisms and generating falsifiable predictions that prospective IRB-approved trials can test.
